% Part: normal-modal-logic
% Chapter: syntax-and-semantics
% Section: modal-validity

\documentclass[../../../include/open-logic-section]{subfiles}

\begin{document}

\olfileid{nml}{syn}{val}

\olsection{માન્યતા}

\begin{explain}
  બધાં નિદર્શનોમાં, એટલે કે દરેક નિદર્શનના દરેક જગતમાં,
  સત્ય હોય તેવાં !!^{formula}s ખાસ રસપ્રદ છે. શક્ય જગતો
  પરના કોઈ સુલભતા સંબંધથી ઉદ્ભવતા અર્થમાં અર્થઘટન
  ``સામાન્ય'' હોય ત્યાં સુધી $\Box$ અને $\Diamond$નું
  અર્થઘટન ગમે તેમ કરવામાં આવે, આ મોડલ વિધાનો
  સત્ય રહે છે. આવાં !!{formula}sને આપણે \emph{માન્ય} કહીએ.
  ઉદાહરણ તરીકે, $\Box(p \land q) \lif \Box p$ માન્ય છે.
  પરંતુ $\Box$ના સત્યલક્ષી અર્થઘટનને આધારે માન્ય હોવાની
  અપેક્ષા રાખી શકાય એવાં કેટલાંક !!{formula}s, જેમ કે
  $\Box p \lif p$, માન્ય નથી. સંબંધાત્મક નિદર્શનો
  રસપ્રદ હોવાનું એક કારણ એ છે કે વિવિધ પ્રકારના
  સુલભતા સંબંધો વડે $\Box$ અને $\Diamond$નાં વિવિધ
  અર્થઘટનો વ્યક્ત કરી શકાય છે. તેથી માન્યતાની વ્યાખ્યા
  માત્ર \emph{બધાં} નિદર્શનો સાપેક્ષે જ નહીં, પરંતુ
  \emph{ચોક્કસ પ્રકારનાં} બધાં નિદર્શનો સાપેક્ષે પણ
  આપવી જોઈએ. દાખલા તરીકે, જે નિદર્શનોમાં દરેક
  જગત પોતામાંથી સુલભ છે, એટલે કે $R$ સ્વવાચક છે,
  તેમાં $\Box p \lif p$ સત્ય છે. નિદર્શનોના વર્ગ
  સાપેક્ષે માન્યતા વ્યાખ્યાયિત કરવાથી આ વાત ટૂંકમાં
  કહી શકાય: $\Box p \lif p$ સ્વવાચક નિદર્શનોના
  વર્ગમાં માન્ય છે.
\end{explain}

\begin{defn}
  !!^a{formula}~$!A$ નિદર્શનોના વર્ગ $\mClass{C}$માં
  \emph{માન્ય} છે, જો તે $\mClass{C}$ના દરેક નિદર્શનમાં
  સત્ય હોય (એટલે કે $\mClass{C}$ના દરેક નિદર્શનના
  દરેક જગતમાં સત્ય હોય). જો $!A$ એ~$\mClass{C}$માં
  માન્ય હોય, તો $\mClass{C} \Entails !A$ લખીએ;
  અને $!A$ \emph{બધાં} નિદર્શનોના વર્ગમાં માન્ય હોય,
  તો $\Entails !A$ લખીએ.
\end{defn}

\begin{prop}\ollabel{prop:subset-class}
  જો $!A$ વર્ગ $\mClass{C}$માં માન્ય હોય, તો તે દરેક
  ઉપવર્ગ $\mClass{C}' \subseteq \mClass{C}$માં પણ માન્ય છે.
\end{prop}

\begin{prop}\ollabel{prop:Nec-rule}
  જો $!A$ માન્ય હોય, તો $\Box!A$ પણ માન્ય છે.
\end{prop}

\begin{proof}
  $\Entails !A$ માનો. $\Entails \Box!A$ બતાવવા માટે
  કોઈપણ નિદર્શન $\mModel{M} =
  \tuple{W, R, V}$ અને $w \in W$ લો. જો $Rww'$ હોય,
  તો $!A$ માન્ય હોવાથી $\mSat{M}{!A}[w']$ છે;
  તેથી $\mSat{M}{\Box!A}[w]$ પણ છે. નિદર્શન
  $\mModel{M}$ અને જગત $w$ ગમે તે હતાં, તેથી
  $\Entails \Box!A$.
\end{proof}

\begin{prob}
  નીચેના સૂત્રો માન્ય છે તે બતાવો:
  \begin{enumerate}
  \item $\Box p \lif \Box (q \lif p)$;
  \item $\Box \lnot \lfalse$;
  \item $\Box p \lif (\Box q \lif \Box p)$.
  \end{enumerate}
\end{prob}

\begin{prob}
  $W = \{w\}$ ધરાવતા નિદર્શનો $\mModel{M} =
  \tuple{W, R, V}$ના વર્ગ $\mClass{C}$માં
  $!A \lif \Box!A$ માન્ય છે તે બતાવો. એ જ રીતે,
  $R = \emptyset$ ધરાવતા નિદર્શનો $\mModel{M} =
  \tuple{W, R, V}$ના વર્ગમાં $!B \lif \Box !A$
  અને $\Diamond !A \lif !B$ માન્ય છે તે બતાવો.
\end{prob}

\end{document}
